\section{Các bên liên quan}

Các bên liên quan chính của hệ thống gồm:

\begin{itemize}
    \item \textbf{Quản trị viên hệ thống}: quản lý cấu hình nền tảng và quyền truy cập khách hàng;
    \item \textbf{Quản trị viên khách hàng}: quản lý người dùng, vai trò, chi nhánh, sản phẩm và nghiệp vụ trong phạm vi tổ chức;
    \item \textbf{Nhân viên kho}: nhập hàng, chuyển kho, kiểm kê và xử lý biến động tồn;
    \item \textbf{Nhân viên bán hàng}: tìm kiếm sản phẩm, giữ hàng và thực hiện thanh toán POS;
    \item \textbf{Quản lý}: theo dõi doanh thu, tồn kho, lợi nhuận và cảnh báo;
    \item \textbf{Kênh bán hàng bên ngoài}: gửi sự kiện đơn hàng đến trung tâm đơn hàng.
\end{itemize}

\section{Yêu cầu chức năng}

\subsection{Xác thực và đa khách hàng}

\begin{itemize}
    \item \textbf{FR-AUTH-01}: xác thực người dùng an toàn;
    \item \textbf{FR-AUTH-02}: cô lập dữ liệu giữa các khách hàng;
    \item \textbf{FR-AUTH-03}: hỗ trợ phân quyền dựa trên vai trò;
    \item \textbf{FR-AUTH-04}: hỗ trợ nhiều chi nhánh trong một khách hàng;
    \item \textbf{FR-AUTH-05}: cấp mã truy cập và mã làm mới cho phiên đăng nhập.
\end{itemize}

\subsection{Danh mục sản phẩm}

\begin{itemize}
    \item \textbf{FR-PROD-01}: quản lý nhóm sản phẩm và thương hiệu;
    \item \textbf{FR-PROD-02}: quản lý sản phẩm, biến thể và SKU;
    \item \textbf{FR-PROD-03}: quản lý thông tin mã vạch;
    \item \textbf{FR-PROD-04}: quản lý giá bán;
    \item \textbf{FR-PROD-05}: tìm kiếm SKU nhanh theo mã, mã vạch hoặc thông tin sản phẩm.
\end{itemize}

\subsection{Quản lý tồn kho}

\begin{itemize}
    \item \textbf{FR-INV-01}: ghi mọi biến động tồn kho vào sổ cái chỉ ghi thêm;
    \item \textbf{FR-INV-02}: tạo phiếu nhập hàng;
    \item \textbf{FR-INV-03}: chuyển hàng giữa các chi nhánh;
    \item \textbf{FR-INV-04}: thực hiện kiểm kê và ghi nhận điều chỉnh đã được phê duyệt;
    \item \textbf{FR-INV-05}: cung cấp lịch sử thay đổi tồn kho có khả năng kiểm toán.
\end{itemize}

\subsection{Tính giá vốn}

Hệ thống phải tính giá vốn bình quân gia quyền (WAC) sau khi tồn kho được bổ sung. Với số
lượng trước đó $Q_{old}$, giá vốn đơn vị $C_{old}$, số lượng nhập thêm $Q_{in}$ và giá nhập
đơn vị $C_{in}$, giá vốn bình quân được tính như sau:

\[
C_{WAC} =
\frac{Q_{old}C_{old} + Q_{in}C_{in}}
     {Q_{old}+Q_{in}}
\]

Với số lượng bán $Q_{sold}$, giá vốn hàng bán được tính:

\[
COGS = Q_{sold} \times C_{WAC}
\]

Giá trị COGS được sử dụng cho đơn hàng đã hoàn tất phải được lưu lại như một ảnh chụp giao dịch,
không phụ thuộc vào các thay đổi giá vốn phát sinh sau đó.

\subsection{Quản lý đơn hàng}

\begin{itemize}
    \item \textbf{FR-ORD-01}: tiếp nhận đơn hàng từ nhiều kênh;
    \item \textbf{FR-ORD-02}: chuẩn hóa dữ liệu riêng của từng kênh thành mô hình đơn hàng chuẩn;
    \item \textbf{FR-ORD-03}: quản lý máy trạng thái đơn hàng;
    \item \textbf{FR-ORD-04}: xử lý xác nhận, hoàn tất và hủy đơn hàng.
\end{itemize}

Vòng đời chính của đơn hàng là:

\[
Draft \rightarrow Reserved \rightarrow Confirmed \rightarrow Completed
\]

Việc hủy được hỗ trợ tại các trạng thái trung gian phù hợp theo quy tắc nghiệp vụ.

\subsection{Giữ hàng và chống bán vượt}

Đối với mỗi SKU và chi nhánh, số lượng có thể bán được xác định bởi:

\[
Available = OnHand - Reserved
\]

Một dòng đơn hàng chỉ được giữ khi số lượng yêu cầu không vượt quá số lượng có thể bán.
Các thao tác giữ, xác nhận, trừ và giải phóng tồn phải được thực hiện trong giao dịch để
ngăn trạng thái tồn kho không nhất quán.

\subsection{Điểm bán hàng POS}

Ứng dụng POS phải:

\begin{itemize}
    \item tìm kiếm sản phẩm và SKU nhanh;
    \item kiểm tra tồn kho trước khi thanh toán;
    \item tạo và hoàn tất giao dịch bán;
    \item phát hành thông tin hóa đơn hoặc biên nhận;
    \item thực hiện chuỗi Giữ $\rightarrow$ Xác nhận $\rightarrow$ Trừ tồn một cách nguyên tử.
\end{itemize}

\subsection{Báo cáo}

Phân hệ báo cáo phải cung cấp:

\begin{itemize}
    \item báo cáo doanh thu và lợi nhuận gộp;
    \item tổng hợp COGS;
    \item báo cáo số dư và biến động tồn kho;
    \item thông tin tốc độ luân chuyển hàng;
    \item cảnh báo tồn kho thấp.
\end{itemize}

\subsection{Mô phỏng và xử lý thời gian thực}

Hệ thống phải hỗ trợ bộ mô phỏng webhook cho sự kiện từ kênh trực tuyến. SignalR được sử dụng
để đẩy thông báo thời gian thực khi phù hợp, còn Hangfire thực hiện các tác vụ nền như đồng bộ
theo hàng đợi và các công việc định kỳ.

\section{Yêu cầu phi chức năng}

\subsection{Hiệu năng}

Các mục tiêu hiệu năng gồm:

\begin{itemize}
    \item tìm kiếm SKU: độ trễ p95 dưới 200 ms;
    \item ghi đơn hàng: độ trễ p95 dưới 500 ms;
    \item báo cáo: dưới 2 giây với tập dữ liệu dưới 100.000 bản ghi;
    \item kịch bản đồng thời: hỗ trợ kiểm thử 50 đơn hàng đồng thời cùng hướng đến một SKU chỉ còn một đơn vị.
\end{itemize}

\subsection{Đồng thời và nhất quán}

Việc triển khai phải sử dụng chiến lược khóa bi quan hoặc khóa lạc quan phù hợp. Các thao tác
tồn kho quan trọng phải được thực hiện nguyên tử và giữ tính nhất quán khi có nhiều yêu cầu
đồng thời.

\subsection{Bảo mật}

Các yêu cầu bảo mật gồm:

\begin{itemize}
    \item băm mật khẩu bằng BCrypt hoặc Argon2;
    \item bảo vệ truyền tải bằng HTTPS/TLS 1.3;
    \item sử dụng mã truy cập JWT và mã làm mới;
    \item phân quyền dựa trên vai trò;
    \item kiểm tra quyền theo khách hàng và giới hạn truy vấn theo khách hàng.
\end{itemize}

\subsection{Khả dụng và khả năng sử dụng}

Mục tiêu khả dụng là 99,5\%. Quy trình POS cho một giao dịch thông thường nên hoàn thành trong
ít hơn ba thao tác chính của người dùng và giao diện phải thích ứng với kích thước màn hình.

\subsection{Khả năng bảo trì và kiểm thử}

Hệ thống phải tuân thủ các nguyên tắc Kiến trúc Sạch. Các mô-đun lõi hướng tới độ bao phủ kiểm
thử đơn vị và tích hợp tối thiểu 80\%, phụ thuộc vào bằng chứng triển khai và kiểm thử cuối cùng.

\section{Ma trận truy vết yêu cầu}

Ma trận truy vết yêu cầu (RTM) ánh xạ từng yêu cầu đến thành phần thiết kế, mô-đun hiện thực
và bằng chứng kiểm thử. Cấu trúc rút gọn được trình bày dưới đây.

\begin{center}
\begin{tabular}{|p{2.5cm}|p{4cm}|p{4cm}|}
\hline
Yêu cầu & Khu vực hiện thực & Xác minh \\
\hline
FR-AUTH & Xác thực, khách hàng, phân quyền & Kiểm thử xác thực và phân quyền \\
\hline
FR-PROD & Phân hệ sản phẩm và SKU & Kiểm thử API và giao diện sản phẩm \\
\hline
FR-INV & Sổ cái tồn kho & Kiểm thử tích hợp tồn kho \\
\hline
FR-COST & Dịch vụ WAC và COGS & Kiểm thử đơn vị tính giá vốn \\
\hline
FR-ORD & Trung tâm đơn hàng & Kiểm thử luồng đơn hàng \\
\hline
FR-RSE & Dịch vụ giữ hàng & Kiểm thử đồng thời \\
\hline
FR-POS & Ứng dụng POS & Kiểm thử thanh toán \\
\hline
FR-REP & Phân hệ báo cáo & Kiểm thử và đối chiếu báo cáo \\
\hline
\end{tabular}
\end{center}
